\chapter{The Short-Time Fourier Transform and Spectrograms}

The Fourier transform is a beautiful machine, but it has a blind spot, and that
blind spot is \emph{time}. Hand it a recording of someone whistling a rising
tone, or a snippet of speech, or the wrist motion of a guitarist strumming, and
it will tell you precisely \emph{which} frequencies are present across the whole
recording. What it will not tell you is \emph{when} each frequency happened. It
flattens the entire history of the signal into one verdict. For a steady,
unchanging signal that verdict is the whole truth. For anything that
\emph{changes} -- which is to say, for almost every interesting signal in the
real world -- it is a smear.

This chapter is about fixing that. The fix is almost embarrassingly simple in
concept: instead of analysing the whole signal at once, chop it into short
pieces and Fourier-analyse each piece. The result is the
\textbf{Short-Time Fourier Transform} (STFT), and when we draw its magnitude as
a picture we get a \textbf{spectrogram} -- arguably the single most useful
visualization in all of signal processing. By the end you will be able to read a
spectrogram the way a musician reads a score, you will understand the iron law
that trades time resolution against frequency resolution, and you will see why
the spectrogram is exactly the right lens for the changing, gesture-driven
signals at the heart of our wrist controller.

\chapterline

\section{Why the plain Fourier transform is not enough}
\label{ch10:why-not-enough}

Recall what the Fourier transform of a signal $x(t)$ does. It computes, for every
frequency $f$, the inner product of $x$ with a complex sinusoid that spins at
that frequency:
\begin{equation}
\label{ch10:ft-defn}
X(f) = \int_{-\infty}^{\infty} x(t)\, e^{-\jj 2\pi f t}\,\dd t
     = \inner{x}{e^{\jj 2\pi f t}}.
\end{equation}
The integral runs over \emph{all} time, from the beginning of the universe to
the end of it. Every instant of the signal contributes to every frequency bin
with equal weight. That is exactly what we want when the signal is
\emph{stationary} -- when its frequency content does not change -- because then
``which frequencies are present'' is a well-posed question with a single answer.

\begin{intuitionbox}
Think of \eqref{ch10:ft-defn} as asking the whole signal one yes/no-ish question
per frequency: ``On the whole, how much do you look like a sine wave wiggling at
$f$ cycles per second?'' If the signal is a pure tone the answer is sharp. But if
the signal was a low tone for the first second and a high tone for the second
second, the transform reports \emph{both} tones, with no hint that they took
turns. It heard the whole concert and told you the set list -- but not the
running order.
\end{intuitionbox}

Here is the trap in concrete form. Consider two signals over $0 \le t < 2$
seconds. Signal A is a 5\,Hz tone for the first second followed by a 15\,Hz tone
for the second second. Signal B is a 15\,Hz tone first, then a 5\,Hz tone. These
are obviously different signals -- if they were sounds you would hear one rising
and the other falling. Yet their magnitude spectra $\abs{X(f)}$ are nearly
identical: both contain ``5\,Hz'' and ``15\,Hz'' in roughly equal amounts. All
the information distinguishing them is hidden in the \emph{phase}, tangled across
the whole frequency axis in a way no human can read.

\begin{trapbox}
``The magnitude spectrum tells me what is in my signal.'' It tells you what is in
your signal \emph{integrated over the entire record}. Two completely different
time-stories can produce the same magnitude spectrum. If your signal evolves --
speech, music, a gesture, an engine spinning up -- the global spectrum is the
wrong tool, and no amount of zero-padding or finer frequency resolution will fix
it. You need to localize in time.
\end{trapbox}

\begin{figure}[H]
\centering
\begin{tikzpicture}
\begin{axis}[
  width=12cm, height=4.2cm,
  xlabel={time $t$ (s)}, ylabel={A($t$)},
  xmin=0, xmax=2, ymin=-1.4, ymax=1.4,
  axis lines=box, samples=200, domain=0:1,
  title={Signal A: 5\,Hz then 15\,Hz},
  every axis title/.style={below right,at={(0,1)}},
]
\addplot[sigblue, thick, domain=0:1, samples=160] {sin(deg(2*pi*5*x))};
\addplot[specred, thick, domain=1:2, samples=180] {sin(deg(2*pi*15*x))};
\draw[gray, dashed] (axis cs:1,-1.4) -- (axis cs:1,1.4);
\end{axis}
\end{tikzpicture}
\caption{A non-stationary signal: a low tone for one second, then a high tone.
Its global magnitude spectrum would show both 5 and 15\,Hz with no clue that they
happened in sequence. Reversing the order leaves the magnitude spectrum almost
unchanged.}
\label{ch10:fig-nonstationary}
\end{figure}

\begin{happeningbox}
What is really happening is a mismatch of \emph{basis} and \emph{signal}. The
Fourier basis $\set{e^{\jj 2\pi f t}}$ consists of functions that are spread out
over all time -- each one is an eternal, never-decaying oscillation. To represent
something local in time (like ``a burst of 15\,Hz around $t = 1.5$\,s'') in such
a basis, you must add up a huge number of these eternal sinusoids whose tails
cancel everywhere except the burst. The information about \emph{when} survives
only in the delicate phase relationships among many bins. The STFT changes the
basis to one made of \emph{localized} pieces, so that ``when'' becomes an axis you
can read directly.
\end{happeningbox}

\section{The idea: look through a moving window}
\label{ch10:moving-window}

The cure follows from the diagnosis. If the problem is that the integral sees all
of time at once, then restrict its view. Before transforming, multiply the signal
by a \textbf{window} $w(t)$ -- a short bump that is near zero everywhere except in
a small neighbourhood. The product $x(t)\,w(t - \tau)$ is a copy of the signal as
seen through a porthole centred at time $\tau$: nonzero only near $\tau$, zero
elsewhere. Fourier-transform \emph{that}, and you learn the frequency content of
the signal \emph{around time $\tau$}. Then slide the window to a new $\tau$ and
repeat.

\begin{definition}[Short-Time Fourier Transform]
\label{ch10:stft-defn}
Given a signal $x(t)$ and a window $w(t)$, the (continuous) STFT is
\begin{equation}
\label{ch10:stft-eqn}
X(\tau, f) \;=\; \int_{-\infty}^{\infty} x(t)\, w(t - \tau)\, e^{-\jj 2\pi f t}\,\dd t .
\end{equation}
It is a function of \emph{two} variables: the time position $\tau$ of the window,
and the frequency $f$. For each fixed $\tau$ it is just an ordinary Fourier
transform of the windowed slice $x(t)\,w(t-\tau)$.
\end{definition}

Compare \eqref{ch10:stft-eqn} with \eqref{ch10:ft-defn}: the only new ingredient
is the factor $w(t-\tau)$. That single factor turns a one-dimensional function of
frequency into a two-dimensional function of time \emph{and} frequency. We have
traded a sharp but timeless answer for a slightly blurry but time-resolved one.

\begin{whybox}
Why does multiplying by a window before transforming localize in time? Because
the integral \eqref{ch10:stft-eqn} can only ``hear'' the signal where the window
is nonzero. Outside the porthole, $w(t-\tau) = 0$ kills the integrand regardless
of what $x$ is doing there. So $X(\tau, f)$ is built \emph{entirely} from the
slice of $x$ near $\tau$. Slide $\tau$ and you get a fresh local spectrum for each
position -- a movie of the spectrum over time, rather than a single snapshot.
\end{whybox}

\begin{intuitionbox}
A spotlight in a dark theatre. The plain Fourier transform turns on the house
lights and asks ``what is on this stage, overall?'' The STFT keeps the house dark
and walks a narrow spotlight ($w$) across the stage, asking at each position
``what is \emph{here}, right now?'' A wide spotlight sees a lot of the stage at
once (poor time localization) but lets your eyes adjust to read fine detail
(good frequency detail). A pinpoint spotlight tells you exactly \emph{where}
something is, but you can barely make out what it is. That tension is the whole
story of this chapter.
\end{intuitionbox}

\subsection{Two ways to read the same equation}

There is a subtlety worth pinning down, because it explains a phase convention you
will meet in real libraries. Equation \eqref{ch10:stft-eqn} multiplies by
$e^{-\jj 2\pi f t}$ using \emph{absolute} time $t$. An equivalent and very common
alternative slides the complex exponential along with the window:
\begin{equation}
\label{ch10:stft-shifted}
\tilde X(\tau, f) \;=\; \int_{-\infty}^{\infty} x(t)\, w(t-\tau)\, e^{-\jj 2\pi f (t-\tau)}\,\dd t
                 \;=\; e^{\jj 2\pi f \tau}\, X(\tau, f).
\end{equation}
The two differ only by the phase factor $e^{\jj 2\pi f \tau}$. Since we will spend
almost all our time looking at the \emph{magnitude} $\abs{X(\tau,f)}$, and
$\abs{e^{\jj 2\pi f \tau}} = 1$, the choice does not affect any spectrogram we
draw. It matters only when you reconstruct the signal or interpret phase. Keep it
in your back pocket and move on.

\begin{connectionbox}
The STFT is, at heart, an \emph{inner product} -- the same notion from the linear
algebra companion. Fix $\tau$ and $f$; then
$X(\tau,f) = \inner{x}{\,g_{\tau,f}}$ where
$g_{\tau,f}(t) = w(t-\tau)\,e^{\jj 2\pi f t}$ is a windowed sinusoid called a
\emph{time-frequency atom} or \emph{Gabor atom}. Each atom is a little wave-packet
parked at time $\tau$ and tuned to frequency $f$. The STFT measures how much of
each atom lives in your signal. This is exactly the projection-onto-a-basis view
of Fourier analysis, now with a basis whose members are local in both time and
frequency. The wavelet transform of Chapter~24 is the same idea with atoms that
stretch and shrink instead of just sliding.
\end{connectionbox}

\section{Windows: the shape of the porthole}
\label{ch10:windows}

We cannot avoid windowing -- even doing nothing is a window. If you simply grab
$N$ consecutive samples and transform them, you have implicitly multiplied the
infinite signal by a \textbf{rectangular window} $\rect$ that is $1$ on those
samples and $0$ elsewhere. The question is never \emph{whether} to window, only
\emph{which shape}.

The shape matters because of what happens at the \emph{edges}. A rectangular
window slams the signal to zero abruptly. In the frequency domain, that abrupt
cut convolves the true spectrum with the transform of the rectangle -- a sinc
function with tall \emph{side lobes} (Chapter~7). Those side lobes leak energy
from a strong tone into neighbouring frequencies, a phenomenon called
\textbf{spectral leakage}. A tapered window that eases the signal gently down to
zero at both ends has much smaller side lobes, at the cost of a slightly wider
central lobe.

\begin{figure}[H]
\centering
\begin{tikzpicture}
\begin{axis}[
  width=12cm, height=4.6cm,
  xlabel={normalized position across window}, ylabel={$w[n]$},
  xmin=0, xmax=1, ymin=-0.05, ymax=1.15,
  axis lines=box, samples=160, domain=0:1,
  legend pos=south east, legend cell align=left,
  title={Three common windows}, every axis title/.style={below right,at={(0,1)}},
]
\addplot[gray, thick] coordinates {(0,1)(1,1)};
\addlegendentry{rectangular}
\addplot[sigblue, thick, domain=0:1, samples=160] {0.5-0.5*cos(deg(2*pi*x))};
\addlegendentry{Hann}
\addplot[specred, thick, domain=0:1, samples=160] {0.54-0.46*cos(deg(2*pi*x))};
\addlegendentry{Hamming}
\end{axis}
\end{tikzpicture}
\caption{The rectangular window has hard edges (large leakage, narrow main lobe).
The Hann and Hamming windows taper smoothly to (or near) zero, trading a wider
main lobe for dramatically lower side lobes.}
\label{ch10:fig-windows}
\end{figure}

The two tapered windows in Figure~\ref{ch10:fig-windows} are workhorses. For a
window of length $N$ (samples $n = 0,\dots,N-1$):
\begin{align}
\text{Hann:}\quad & w[n] = \tfrac12\!\left(1 - \cos\tfrac{2\pi n}{N-1}\right), \label{ch10:hann}\\
\text{Hamming:}\quad & w[n] = 0.54 - 0.46\cos\tfrac{2\pi n}{N-1}. \label{ch10:hamming}
\end{align}
The Hann window goes all the way to zero at both ends; the Hamming window stops
just shy of zero, which cancels the nearest side lobe even better but leaves a
small pedestal far out. For everyday spectrograms the Hann window is a fine
default and the one we will assume unless noted.

\begin{focusbox}
A window does two jobs at once, and it is worth holding both in mind:
\begin{enumerate}
  \item \textbf{It localizes in time.} The window's \emph{width} sets how big a
        slice of signal each spectrum sees.
  \item \textbf{It shapes the frequency response.} The window's \emph{taper} sets
        how cleanly that slice's spectrum is rendered -- how much leakage there is.
\end{enumerate}
Width controls the time/frequency tradeoff (next section); taper controls leakage
within whatever resolution the width allows.
\end{focusbox}

\begin{connectionbox}
Multiplication in time is \emph{convolution} in frequency -- the convolution
theorem of Chapter~11. The windowed slice $x(t)w(t-\tau)$ therefore has spectrum
$X(f) \conv W(f)$ (up to a phase from the shift). Each true spectral line gets
replaced by a blurred copy of the window's own spectrum $W(f)$. Choosing a window
is choosing the shape of that blur. This is the same machinery that governs FIR
filter design in Chapter~16, where windows tame the ripples of a truncated ideal
filter.
\end{connectionbox}

\section{The spectrogram: a picture of the STFT}
\label{ch10:spectrogram}

The STFT $X(\tau, f)$ is a complex-valued function of two real variables. We
cannot easily look at a complex 2-D surface, so we throw away the phase and look
at the squared magnitude. That image is the spectrogram.

\begin{definition}[Spectrogram]
\label{ch10:specgram-defn}
The spectrogram of $x$ is the squared magnitude of its STFT,
\begin{equation}
\label{ch10:specgram-eqn}
S(\tau, f) \;=\; \abs{X(\tau, f)}^2,
\end{equation}
displayed as an image with time $\tau$ on the horizontal axis, frequency $f$ on
the vertical axis, and $S$ encoded as brightness or colour. Because spectral
energy spans many orders of magnitude, it is almost always shown in decibels,
$10\log_{10} S$.
\end{definition}

\begin{focusbox}
\textbf{How to read a spectrogram.}
\begin{itemize}
  \item \textbf{Horizontal axis = time.} Left is earlier, right is later. Read it
        like a piano roll or a music score scrolling past.
  \item \textbf{Vertical axis = frequency.} Bottom is low (bass, slow), top is
        high (treble, fast).
  \item \textbf{Colour/brightness = energy.} Bright means strong; dark means
        little or no energy at that time and frequency.
  \item \textbf{Horizontal bright lines} = sustained tones (a held note, a hum, a
        steady rotation).
  \item \textbf{Vertical bright streaks} = sudden broadband events (a click, a
        drum hit, a footstep, a tap of the finger).
  \item \textbf{Sloping lines} = glides -- a rising line is a frequency sweeping
        up, like a whistle going up or a motor spinning faster.
  \item \textbf{Stacks of evenly spaced lines} = a harmonic tone (a voiced
        vowel, a plucked string); the spacing is the fundamental.
\end{itemize}
\end{focusbox}

The schematic in Figure~\ref{ch10:fig-specgram} shows the spectrogram of our
running example -- the rising signal of Figure~\ref{ch10:fig-nonstationary} --
with the time-ordering now plainly visible. The plain Fourier transform could
never show this.

\begin{figure}[H]
\centering
\begin{tikzpicture}
\begin{axis}[
  width=12cm, height=5.2cm,
  xlabel={time $\tau$ (s)}, ylabel={frequency $f$ (Hz)},
  xmin=0, xmax=2, ymin=0, ymax=20,
  axis lines=box,
  title={Spectrogram of ``5\,Hz then 15\,Hz'' (schematic)},
  every axis title/.style={below right,at={(0,1)}},
]
\addplot[draw=none, fill=specred!75] coordinates {(0.02,4.3)(1.0,4.3)(1.0,5.7)(0.02,5.7)} \closedcycle;
\addplot[draw=none, fill=specred!75] coordinates {(1.0,14.3)(1.98,14.3)(1.98,15.7)(1.0,15.7)} \closedcycle;
\draw[gray, dashed] (axis cs:1,0) -- (axis cs:1,20);
\node at (axis cs:0.5,9) {\small 5\,Hz};
\node at (axis cs:1.5,11) {\small 15\,Hz};
\end{axis}
\end{tikzpicture}
\caption{The spectrogram of the signal in Figure~\ref{ch10:fig-nonstationary}.
The low tone occupies the first second, the high tone the second second --
exactly the ``running order'' the global Fourier transform discarded. Each bright
patch has nonzero height (frequency spread) and width (time spread) set by the
window; an ideal point would require a window that is both infinitely short and
infinitely long.}
\label{ch10:fig-specgram}
\end{figure}

\begin{applicationbox}
\textbf{The voiceprint.} The spectrogram was born in wartime speech research at
Bell Labs in the 1940s (the ``sound spectrograph''), and it remains the native
representation of speech. A vowel shows up as a stack of harmonics shaped by
broad bright bands called \emph{formants} -- resonances of the vocal tract -- that
slide as the mouth changes shape. Consonants such as /s/ and /f/ appear as
high-frequency noise fills; plosives like /p/ and /t/ as a brief silence then a
vertical burst. Trained readers can literally read speech off a spectrogram. Every
modern speech recognizer still starts by computing something very close to a
spectrogram (the mel-spectrogram of Chapter~27).
\end{applicationbox}

\begin{historybox}
The sound spectrograph was developed at Bell Telephone Laboratories during World
War II and declassified afterward; the 1947 book \emph{Visible Speech} by Potter,
Kopp, and Green popularized reading speech from these images and even explored
teaching the deaf to speak by watching their own spectrograms in near real time.
The mathematics was older -- Dennis Gabor's 1946 paper \emph{Theory of
Communication} introduced the windowed transform with a Gaussian window and the
``logon'' (what we now call a time-frequency atom), and proved the uncertainty
limit we meet in the next section. The spectrogram is thus the visible face of
Gabor's analysis.
\end{historybox}

\section{The central tradeoff: time versus frequency resolution}
\label{ch10:tradeoff}

Here is the iron law of the STFT, and it is not a defect of any particular
algorithm -- it is a law of nature about waves. \emph{You cannot be arbitrarily
precise in time and in frequency at the same time.} The window length forces a
choice.

Think about it physically. To measure a frequency precisely you must watch an
oscillation for many cycles -- a single glimpse cannot tell a 100\,Hz tone from a
101\,Hz tone, because in a few milliseconds they look identical. Resolving 1\,Hz
of difference requires watching for about 1 second, because only then have the two
tones drifted a full cycle out of step. But watching for a full second means you
have \emph{no idea when} within that second the tone occurred -- you have blurred
time by a second to sharpen frequency to a hertz.

\begin{whybox}
Why is this unavoidable? Because frequency \emph{is} a statement about behaviour
over an interval of time. ``This signal has frequency $f$'' is shorthand for
``this signal repeats every $1/f$ seconds, for a while.'' If you only look for an
instant, the phrase ``for a while'' has no meaning, so neither does a precise
frequency. The shorter your window, the fewer cycles fit inside it, and the
fuzzier the frequency you can name. Conversely, the longer the window, the more
time you have smeared together, and the vaguer your sense of \emph{when}. The two
resolutions are reciprocals of one window length; tightening one necessarily
loosens the other.
\end{whybox}

We can make this quantitative. If the window spans a duration $T$ (in seconds),
then:
\begin{equation}
\label{ch10:res-tradeoff}
\Delta t \approx T \qquad\text{and}\qquad \Delta f \approx \frac{1}{T},
\qquad\text{so}\qquad \Delta t \cdot \Delta f \approx 1 .
\end{equation}
The product of the time spread and the frequency spread is roughly constant. You
can shuttle resolution back and forth between the axes by changing $T$, but you
cannot shrink the product below a fixed floor. Figure~\ref{ch10:fig-tiling} draws
this as a tiling of the time-frequency plane: each window length reshapes the
tiles but holds their area fixed. With a Gaussian window the floor is
exactly $\Delta t\,\Delta f \ge \frac{1}{4\pi}$ -- this is the
\emph{Gabor--Heisenberg uncertainty} relation, the same inequality (up to
constants) that governs position and momentum in quantum mechanics, for the same
mathematical reason.

\begin{formalbox}
Define the spreads of a unit-energy window $w$ by its second moments about its
centre,
\[
(\Delta t)^2 = \int t^2\,\abs{w(t)}^2\,\dd t, \qquad
(\Delta f)^2 = \int f^2\,\abs{W(f)}^2\,\dd f .
\]
Then for \emph{any} window,
\[
\Delta t \cdot \Delta f \;\ge\; \frac{1}{4\pi},
\]
with equality \emph{only} for a Gaussian $w(t) = e^{-a t^2}$. No window beats the
Gaussian; no window escapes the bound. This is why the Gaussian-windowed STFT (the
\emph{Gabor transform}) is the theoretically optimal trade and why ``you can't
have both'' is a theorem, not an excuse.
\end{formalbox}

\begin{figure}[H]
\centering
\begin{tikzpicture}[scale=1.0]
% short window: tall thin tiles
\begin{scope}
  \draw[->,gray] (0,0) -- (3.6,0) node[right,font=\small]{time};
  \draw[->,gray] (0,0) -- (0,3.4) node[above,font=\small]{freq};
  \foreach \x in {0,0.6,1.2,1.8,2.4}{
    \foreach \y in {0,1.0,2.0}{
      \draw[sigblue!70,fill=sigblue!8] (\x,\y) rectangle (\x+0.6,\y+1.0);
    }}
  \node[font=\small] at (1.5,-0.55) {short window};
\end{scope}
% long window: short wide tiles
\begin{scope}[xshift=5.2cm]
  \draw[->,gray] (0,0) -- (3.6,0) node[right,font=\small]{time};
  \draw[->,gray] (0,0) -- (0,3.4) node[above,font=\small]{freq};
  \foreach \x in {0,1.5,3.0}{
    \foreach \y in {0,0.4,0.8,1.2,1.6,2.0,2.4}{
      \draw[specred!70,fill=specred!8] (\x,\y) rectangle (\x+1.5,\y+0.4);
    }}
  \node[font=\small] at (1.5,-0.55) {long window};
\end{scope}
\end{tikzpicture}
\caption{The time-frequency plane tiled by STFT atoms. Every tile has the same
\emph{area} ($\Delta t\,\Delta f \approx 1$), so a short window gives tall, thin
tiles (good time resolution, coarse frequency) and a long window gives short, wide
tiles (good frequency resolution, coarse time). Changing the window length reshapes
the tiles but never shrinks their area.}
\label{ch10:fig-tiling}
\end{figure}

\subsection{Wideband versus narrowband}

The same signal looks qualitatively different depending on which way you spend the
tradeoff. The two regimes even have names borrowed from speech analysis:

\begin{itemize}
  \item \textbf{Narrowband spectrogram} (long window, e.g.\ 30--50\,ms of speech):
        fine \emph{frequency} resolution, coarse time. Individual harmonics of a
        voiced sound resolve into separate horizontal lines, but fast events
        smear. Good for ``what notes / which pitch.''
  \item \textbf{Wideband spectrogram} (short window, e.g.\ 3--6\,ms of speech):
        fine \emph{time} resolution, coarse frequency. Harmonics blur together
        into broad formant bands, but each glottal pulse shows as a crisp vertical
        striation. Good for ``what is happening when / timing of events.''
\end{itemize}

Figure~\ref{ch10:fig-tradeoff} contrasts the two for a single tone that jumps in
frequency. Neither view is ``correct''; they answer different questions.

\begin{figure}[H]
\centering
\begin{tikzpicture}
\begin{axis}[
  name=left, width=6.4cm, height=4.6cm,
  xlabel={time}, ylabel={freq},
  xmin=0, xmax=2, ymin=0, ymax=20, axis lines=box,
  title={short window: sharp in time}, every axis title/.style={below right,at={(0,1)},font=\small},
  xtick=\empty, ytick=\empty,
]
\addplot[draw=none, fill=specred!30] coordinates {(0.04,1)(1.0,1)(1.0,9)(0.04,9)} \closedcycle;
\addplot[draw=none, fill=specred!30] coordinates {(1.0,11)(1.96,11)(1.96,19)(1.0,19)} \closedcycle;
\draw[gray,dashed] (axis cs:1,0)--(axis cs:1,20);
\end{axis}
\begin{axis}[
  name=right, at={(left.east)}, xshift=1.1cm, anchor=west,
  width=6.4cm, height=4.6cm,
  xlabel={time}, ylabel={freq},
  xmin=0, xmax=2, ymin=0, ymax=20, axis lines=box,
  title={long window: sharp in freq}, every axis title/.style={below right,at={(0,1)},font=\small},
  xtick=\empty, ytick=\empty,
]
\addplot[draw=none, fill=specred!70] coordinates {(0.25,4.4)(1.15,4.4)(1.15,5.6)(0.25,5.6)} \closedcycle;
\addplot[draw=none, fill=specred!70] coordinates {(0.85,14.4)(1.75,14.4)(1.75,15.6)(0.85,15.6)} \closedcycle;
\draw[gray,dashed] (axis cs:1,0)--(axis cs:1,20);
\end{axis}
\end{tikzpicture}
\caption{The same tone-jump (low then high), analysed two ways. Short window
(left): the jump time is crisp but each tone is a fat vertical smear in frequency.
Long window (right): each tone is a thin horizontal band in frequency, but the
transition time is blurred and the two patches even overlap in time. The
window length picks your poison.}
\label{ch10:fig-tradeoff}
\end{figure}

\begin{trapbox}
``I will just zero-pad my short window to get more frequency bins.'' Zero-padding
gives you more closely spaced \emph{samples} of the spectrum -- a smoother-looking
curve -- but it does \emph{not} improve frequency \emph{resolution}. Two tones
that were unresolvable in a 6\,ms window stay unresolvable no matter how many
zeros you append; you are interpolating, not informing. True resolution comes only
from a longer \emph{real} window, i.e.\ more genuine signal. Resolution is set by
how long you looked, not by how finely you grid the answer.
\end{trapbox}

\begin{walkbox}
You want to resolve two musical notes a semitone apart near middle A (440\,Hz and
about 466\,Hz, a gap of $\approx 26$\,Hz). What window length do you need? From
\eqref{ch10:res-tradeoff}, $\Delta f \approx 1/T$, so to get $\Delta f \le 26$\,Hz
you need $T \gtrsim 1/26 \approx 38$\,ms. At a 44.1\,kHz sample rate that is about
$0.038 \times 44100 \approx 1700$ samples, so a window of $N = 2048$ is a sensible
choice. Now ask the opposite: with a 38\,ms window, how well can you time a note
onset? Only to about $\pm 38$\,ms -- noticeably loose for fast passages. That is
the tradeoff in a single concrete breath.
\end{walkbox}

\section{From integral to algorithm: the discrete STFT}
\label{ch10:discrete}

Real signals are sampled, and real spectra are computed with the FFT
(Chapter~9). The discrete STFT replaces the integral of
\eqref{ch10:stft-eqn} with a sum over a length-$N$ window, hops the window forward
by $H$ samples (the \textbf{hop size}), and takes an $N$-point DFT of each frame:
\begin{equation}
\label{ch10:dstft}
X[m, k] \;=\; \sum_{n=0}^{N-1} x[\,mH + n\,]\; w[n]\; e^{-\jj 2\pi k n / N},
\qquad k = 0,\dots,N-1 .
\end{equation}
Here $m$ indexes the \emph{frame} (window position, $\tau = mH/f_s$) and $k$
indexes the \emph{frequency bin} ($f = k f_s / N$, where $f_s$ is the sample
rate). Each column of the spectrogram is one frame's FFT; you build the image
column by column as you slide along.

\begin{focusbox}
The four knobs of a discrete STFT, and what each one buys you:
\begin{center}
\begin{tabular}{@{}lll@{}}
\toprule
Knob & Symbol & Controls \\
\midrule
window length & $N$ & time/frequency resolution tradeoff \\
window shape & $w[n]$ & leakage / side-lobe level \\
hop size & $H$ & overlap, smoothness, redundancy, cost \\
FFT length & $\ge N$ & bin spacing (display, not resolution) \\
\bottomrule
\end{tabular}
\end{center}
\end{focusbox}

\subsection{Overlap and the hop size}

If you set the hop equal to the window length ($H = N$), the frames sit side by
side with no overlap. That is cheap but risky: a tapered window suppresses the
signal near each frame's edges, so an event landing on a window boundary is
\emph{seen weakly by both neighbours} and can nearly vanish. The cure is
\textbf{overlap}: advance the window by less than its length, so successive frames
share samples. A hop of $H = N/2$ gives 50\% overlap; $H = N/4$ gives 75\%.
Overlap makes the spectrogram smoother in time, ensures every sample is covered
where some window is near its peak, and is required for clean reconstruction.

\begin{figure}[H]
\centering
\begin{tikzpicture}[scale=1.0]
\def\W{3.2}
\foreach \s/\c in {0/sigblue, 1.6/specred, 3.2/targetgreen}{
  \begin{scope}[xshift=\s cm]
    \draw[\c, thick, domain=0:\W, samples=60, smooth, variable=\t]
      plot ({\t},{0.5-0.5*cos(deg(2*pi*\t/\W))});
  \end{scope}
}
\draw[->,gray] (-0.3,0) -- (7.3,0) node[right]{\small samples};
\draw[->,gray] (-0.3,0) -- (-0.3,1.25) node[above]{\small $w$};
\draw[<->] (0,-0.35) -- (3.2,-0.35) node[midway,below,font=\small]{window $N$};
\draw[<->] (0,-0.95) -- (1.6,-0.95) node[midway,below,font=\small]{hop $H=N/2$};
\end{tikzpicture}
\caption{Three Hann windows at 50\% overlap ($H = N/2$). Where one window dips,
its neighbour is rising, so every sample is well covered. With careful window/hop
choices the overlapping windows even sum to a constant -- the
constant-overlap-add property that lets you reconstruct the signal exactly.}
\label{ch10:fig-overlap}
\end{figure}

\begin{whybox}
Why overlap rather than just a longer window? Overlap does \emph{not} change the
time/frequency resolution -- that is still set by $N$ alone. What overlap buys is
\emph{temporal sampling} of the time-frequency plane. Think of each frame as a
snapshot; with no overlap you photograph the scene every $N$ samples and may miss
fast action between shots. With overlap you photograph more often, catching
transitions and giving a fluid, alias-free movie of the spectrum over time. It
costs more computation (more FFTs) but no resolution penalty.
\end{whybox}

\begin{labbox}
Compute a spectrogram yourself with a few lines of NumPy and watch the knobs bite.
Try a chirp first -- a tone whose frequency rises linearly -- because its
spectrogram should be a clean diagonal line, and any blur or staircase you see is
the tradeoff made visible.
\begin{lstlisting}[language=Python]
import numpy as np
import matplotlib.pyplot as plt

fs = 8000                       # sample rate (Hz)
t  = np.arange(0, 2.0, 1/fs)    # 2 seconds
# linear chirp: instantaneous freq goes 200 -> 1800 Hz
x  = np.sin(2*np.pi*(200*t + (800/2)*t**2))

N, H = 256, 64                  # window length, hop (75% overlap)
w    = np.hanning(N)
frames = range(0, len(x)-N, H)
S = np.array([np.abs(np.fft.rfft(x[i:i+N]*w))**2 for i in frames]).T
S_db = 10*np.log10(S + 1e-12)

plt.imshow(S_db, origin='lower', aspect='auto',
           extent=[0, 2, 0, fs/2])
plt.xlabel('time (s)'); plt.ylabel('frequency (Hz)')
plt.show()
# Now rerun with N=1024, H=256. The diagonal gets thinner
# vertically (better freq) but you see less precisely *where*
# along it the energy is -- the tradeoff, in one experiment.
\end{lstlisting}
\end{labbox}

\subsection{Inverting the STFT: overlap-add}

Because the STFT is redundant (with overlap we store more numbers than we started
with), we can usually throw the signal away and rebuild it from $X[m,k]$. The
recipe is \textbf{overlap-add}: inverse-FFT each frame back to a short time slice,
multiply by a synthesis window, and add the slices back into place with the same
hop. If the analysis and synthesis windows satisfy the \emph{constant
overlap-add} (COLA) condition -- their shifted-by-$H$ copies sum to a constant --
the reconstruction is exact. This invertibility is what makes the STFT a
processing tool, not just a viewer: modify $X[m,k]$ (filter, denoise,
pitch-shift, time-stretch) and resynthesize. It is the engine behind the phase
vocoder and countless audio effects in Chapter~28.

\section{Worked example: reading numbers off the discrete STFT}
\label{ch10:worked}

\begin{example}[Grids and resolution]
\label{ch10:ex-numbers}
Let $f_s = 1000$\,Hz and take a window of $N = 200$ samples with hop $H = 100$
(50\% overlap). What are the time and frequency grids, and what resolution do we
get?

\textbf{Frequency grid.} Bin $k$ sits at $f_k = k f_s/N = k\cdot 1000/200 = 5k$
Hz. So bins are spaced 5\,Hz apart: $0, 5, 10, \dots$ up to the Nyquist frequency
of 500\,Hz at $k = 100$.

\textbf{Frequency resolution.} The window spans $T = N/f_s = 200/1000 = 0.2$\,s,
so by \eqref{ch10:res-tradeoff} $\Delta f \approx 1/T = 5$\,Hz. (With a Hann
window the true main-lobe width is roughly $2/T = 10$\,Hz, a bit coarser than the
bin spacing -- a reminder that bin spacing and resolution are not the same thing.)

\textbf{Time grid.} Frame $m$ is centred near $\tau_m = (mH + N/2)/f_s
= (100m + 100)/1000$\,s, i.e.\ frames step every $H/f_s = 0.1$\,s.

\textbf{Time resolution.} Each frame blurs time by about $T = 0.2$\,s. So an event
shorter than 0.2\,s cannot be timed more precisely than that, even though the
frame \emph{spacing} is a finer 0.1\,s -- again, sampling density (0.1\,s) is not
resolution (0.2\,s). To time events to 20\,ms you would shorten the window to
$N \approx 20$, accepting a coarse $\Delta f \approx 50$\,Hz in return.
\end{example}

\section{Applications: speech, music, motion}
\label{ch10:applications}

\begin{applicationbox}
\textbf{Music transcription and the piano roll.} A spectrogram of a melody is
already most of the way to sheet music: each note is a horizontal segment whose
height is its pitch and whose length is its duration; chords are vertical stacks;
vibrato is a wavy line. Automatic music transcription, the Shazam audio
fingerprint (which hashes constellations of spectrogram peaks), and
pitch-correction plug-ins all live in the STFT domain. A long window resolves
close pitches; a short window pins down fast rhythmic onsets -- so transcribers
often run \emph{both} and combine.
\end{applicationbox}

\begin{applicationbox}
\textbf{Machinery and biomedical monitoring.} A bearing beginning to fail injects
a rising whine and periodic clicks into a machine's vibration; on a spectrogram
this is a new sloping line plus regularly spaced vertical streaks, visible long
before a human would hear it. In an EEG, a seizure announces itself as a sudden
broadband brightening; sleep stages have characteristic spectral signatures
(spindles, slow waves) that clinicians read straight off the spectrogram. The
common thread: \emph{the interesting information is in how the spectrum changes},
which is precisely what the spectrogram makes visible.
\end{applicationbox}

\section{Toward the wrist tracker}
\label{ch10:wrist}

\begin{goalbox}
Gesture signals are \emph{non-stationary by definition} -- a gesture is a change
over time, so a single global spectrum of a recording session is almost useless.
This is exactly the setting the STFT was built for, and the spectrogram is our
primary lens on wrist-worn sensor data.

\textbf{IMU (accelerometer + gyroscope).} Picture the forearm during a tremolo
strum or a rhythmic tap. The accelerometer and gyro signals are oscillatory with
a frequency set by the motion's tempo, riding on slowly drifting baselines from
gravity and arm reorientation. A spectrogram of the gyro magnitude turns a strum
at 6 strokes per second into a bright horizontal line at $\approx 6$\,Hz that
\emph{appears when the strumming starts and vanishes when it stops}, and a single
sharp tap into a clean vertical broadband streak (Figure~\ref{ch10:fig-imu}).
Reading tempo, detecting onset
and offset, and separating ``holding still'' from ``moving rhythmically'' all
become trivial in this view, when they were tangled in the raw trace.

\textbf{Sensible numbers.} Hand/wrist motion energy lives roughly in 0--20\,Hz
(voluntary motion is mostly below 5\,Hz; physiological tremor sits around
8--12\,Hz). Sample the IMU at $f_s = 100$--200\,Hz so Nyquist (Chapter~18)
comfortably clears the band. A window of $T \approx 0.5$\,s ($N = 64$ at 128\,Hz)
gives $\Delta f \approx 2$\,Hz -- enough to tell tremor from voluntary tapping --
while still updating fast enough to feel responsive. Use 50--75\% overlap so the
controller reacts within one hop, $\sim 100$\,ms, instead of one whole window.

\textbf{The limits, stated plainly.} A wrist-mounted IMU senses the motion of the
\emph{forearm}, full stop. Its spectrogram can reveal \emph{coarse, discrete}
gestures -- rhythmic versus still, the tempo of a repetitive motion, a flick, a
twist, a tap -- because those are gross forearm dynamics. It \emph{cannot} recover
which individual finger moved or by how much: finger flexions barely perturb the
wrist's inertial signal, and there is no frequency band that magically encodes
finger identity. For finger-level events you need a sensor that listens to the
muscles or tendons -- surface EMG or force myography on the forearm
(Chapter~37) -- and for continuous finger joint angles you need sensors
\emph{on the fingers} (instrumented rings or gloves). The STFT is a superb tool
for what the IMU can actually see; it cannot conjure information the sensor never
captured.
\end{goalbox}

\begin{figure}[H]
\centering
\begin{tikzpicture}
\begin{axis}[
  width=12cm, height=4.2cm,
  xlabel={time $t$ (s)}, ylabel={gyro mag.},
  xmin=0, xmax=2, ymin=-1.5, ymax=1.7,
  axis lines=box,
  title={Forearm gyro: rhythmic strum, then a sharp tap},
  every axis title/.style={below right,at={(0,1)}},
]
% ~6 Hz oscillation for the strum
\addplot[sigblue, thick, domain=0:1.3, samples=160] {0.8*sin(deg(2*pi*6*x))};
% quiet between
\addplot[gray, thick, domain=1.3:1.55, samples=20] {0};
% sharp tap: a brief spike
\addplot[specred, thick, domain=1.55:1.7, samples=40] {1.4*exp(-((x-1.6)*60)^2)};
\addplot[gray, thick, domain=1.7:2, samples=20] {0};
\node[font=\small, sigblue] at (axis cs:0.65,1.35) {strum $\approx$ 6 Hz};
\node[font=\small, specred] at (axis cs:1.78,1.45) {tap};
\end{axis}
\end{tikzpicture}
\caption{A schematic forearm-gyro trace. The rhythmic strum is a sustained
oscillation (a horizontal line on the spectrogram near 6\,Hz); the tap is a brief
impulse (a vertical broadband streak). The spectrogram separates these two events
that share the same raw amplitude range.}
\label{ch10:fig-imu}
\end{figure}

\begin{buildbox}
\textbf{A streaming spectrogram on a microcontroller.} You do not need to store
the whole recording -- the STFT is naturally incremental, one frame at a time.
Keep a circular buffer of the last $N$ IMU samples; every $H$ samples, copy the
buffer, multiply by a precomputed Hann window, run a fixed-point FFT (CMSIS-DSP's
\texttt{arm\_rfft\_q15} on a Cortex-M, or KissFFT on an ESP32), and you have one
spectrogram column.
\begin{lstlisting}[language=C]
// fires every H new samples (e.g. H=32 at fs=128 Hz)
void on_hop(int16_t *ring, int N) {
    int16_t frame[N];
    for (int n = 0; n < N; n++)
        frame[n] = mul_q15(ring[n], hann_lut[n]); // windowing
    arm_rfft_q15(&S, frame, spectrum);            // in-place FFT
    for (int k = 0; k < N/2; k++)
        mag[k] = mag_q15(spectrum[2*k], spectrum[2*k+1]);
    classify(mag);   // is there 6 Hz energy? a broadband spike?
}
\end{lstlisting}
With $N = 64$ this is a few thousand multiply-accumulates per hop -- trivial even
for an 80\,MHz part -- so a battery wrist unit can stream spectrogram features in
real time within a tight power budget. Precompute the Hann window into a lookup
table; never call \texttt{cos} in the hot loop.
\end{buildbox}

\begin{aestheticbox}
The deep idea of this chapter is that \emph{representation is everything}. The raw
IMU trace and its spectrogram contain the same numbers, related by a linear,
invertible transform -- no information is created. Yet ``is the wrist strumming,
and how fast?'' is nearly invisible in one representation and screaming obvious in
the other. Good signal processing is, over and over, the art of rotating your data
into the basis where the answer is easy to see. The STFT rotates a stream of
samples into a time-frequency picture, and suddenly motion, music, and speech all
become things you can simply \emph{look at}.
\end{aestheticbox}

\section{Summary}
\label{ch10:summary}

\begin{itemize}
  \item The global Fourier transform tells you \emph{what} frequencies are present
        but not \emph{when}; two signals with different time-orderings can share a
        magnitude spectrum.
  \item The \textbf{STFT} fixes this by transforming short windowed slices,
        producing $X(\tau,f)$, a function of both time and frequency
        (Def.~\ref{ch10:stft-defn}). Each value is an inner product of the signal
        with a localized time-frequency atom.
  \item The \textbf{spectrogram} $\abs{X(\tau,f)}^2$ (usually in dB) is the image
        of the STFT: time across, frequency up, energy as brightness. Horizontal
        lines are tones, vertical streaks are transients, slopes are glides.
  \item \textbf{Windows} (Hann, Hamming, \dots) localize in time and control
        spectral leakage; the rectangular window is the no-taper default with the
        worst leakage.
  \item The \textbf{time/frequency tradeoff} is a law: $\Delta t\,\Delta f
        \approx 1$. Long windows resolve frequency (narrowband); short windows
        resolve time (wideband). Zero-padding refines display, not resolution.
  \item \textbf{Overlap} (small hop $H$) samples the time-frequency plane more
        densely and enables exact overlap-add reconstruction, without changing
        resolution.
  \item For the wrist controller, the spectrogram of IMU data exposes rhythmic
        motion, tempo, onsets, and tremor -- but an IMU alone sees only coarse
        forearm motion, not individual fingers.
\end{itemize}

\section{Exercises}
\label{ch10:exercises}

\begin{exercise}
\label{ch10:ex-order}
Two signals last 2\,s: signal A is 100\,Hz then 300\,Hz (1\,s each); signal B is
300\,Hz then 100\,Hz. Argue why their global magnitude spectra are nearly
identical, and sketch the two spectrograms to show how they differ. What feature
of the spectrogram encodes the ordering?
\end{exercise}

\begin{exercise}
\label{ch10:ex-resolution}
You sample at $f_s = 2000$\,Hz and want to distinguish two tones 8\,Hz apart.
(a)~What minimum window duration $T$ do you need? (b)~How many samples $N$ is
that? (c)~With that window, to what precision (in ms) can you locate a sudden
onset? Comment on whether this window is good for timing fast events.
\end{exercise}

\begin{exercise}
\label{ch10:ex-zeropad}
A colleague computes a 64-sample STFT frame, zero-pads it to 1024 points, and
claims the resulting smooth spectrum gives ``16 times better frequency
resolution.'' Explain precisely what zero-padding does and does not improve, and
state what they would actually need to change to resolve closer tones.
\end{exercise}

\begin{exercise}
\label{ch10:ex-hop}
With window length $N = 512$ at $f_s = 16$\,kHz, compute the time between frames
for hops of (a)~$H = 512$, (b)~$H = 256$, (c)~$H = 128$. For each, give the
overlap percentage and the spectrogram's frame rate (columns per second). Which
would you choose to capture drum onsets, and why?
\end{exercise}

\begin{exercise}
\label{ch10:ex-window-leak}
Consider a single 50\,Hz tone sampled at 1000\,Hz over a 256-sample window.
Describe how the spectrogram column differs when you use a rectangular versus a
Hann window. Which shows more leakage into neighbouring bins, and which has the
wider main lobe? State the tradeoff in one sentence.
\end{exercise}

\begin{exercise}
\label{ch10:ex-chirp}
A linear chirp sweeps from 100\,Hz to 900\,Hz over 1\,s. (a)~Describe the ideal
spectrogram. (b)~If you use a very long window, what artifact appears along the
diagonal and why? (c)~If you use a very short window, what artifact appears
instead? Relate both to $\Delta t\,\Delta f \approx 1$.
\end{exercise}

\begin{exercise}
\label{ch10:ex-imu}
A forearm IMU is sampled at 128\,Hz while the user strums rhythmically at about
5 strokes per second, occasionally tapping the table sharply. (a)~On a
spectrogram of the gyro magnitude, what does the steady strumming look like?
(b)~What does a single tap look like? (c)~Choose a window length $N$ (a power of
two) that resolves the 5\,Hz strumming from a 10\,Hz tremor, and justify it with
the resolution formula.
\end{exercise}

\begin{exercise}
\label{ch10:ex-claims}
Your teammate proposes using the wrist-IMU spectrogram to detect \emph{which}
finger the user pressed, reasoning that ``each finger must vibrate at its own
frequency.'' Explain, citing the physics of what an IMU measures, why this will
not work, and name the sensor modalities (and where they must be placed) that
\emph{can} recover finger-level information.
\end{exercise}

\begin{exercise}
\label{ch10:ex-cola}
Explain the constant-overlap-add (COLA) condition in your own words. Why is it
required for exact reconstruction from a modified STFT? Give one example of an
audio effect that is implemented by modifying STFT frames and resynthesizing with
overlap-add.
\end{exercise}

\begin{exercise}
\label{ch10:ex-gabor}
The Gabor (Gaussian-windowed) STFT achieves the theoretical minimum
$\Delta t\,\Delta f = 1/(4\pi)$. (a)~Why might you still prefer a Hann window in
practice, even though it is not optimal? (b)~In what sense is ``optimal'' here a
statement about a \emph{single} atom rather than about reading a whole
spectrogram?
\end{exercise}
